% =====================================================================
%  Convergence of EF21-MuonUSign and EF21-MuonSign by exact reduction to
%  the EF21-Muon framework of Gruntkowska, Burlachenko & Richtarik (2025),
%  arXiv:2510.00643 (source: EF_Muon_Friends/template.tex).
%
%  We do NOT restate their (very long) per-layer bounds. The reduction rests
%  on one non-trivial fact -- the scaled sign is a contractive compressor
%  (Lemma) -- plus routine bookkeeping (LMO step, momentum, loop order). The
%  rates and stepsize rules are stated in our notation; the exact
%  non-asymptotic constants are their Theorems 19 and 24 under the
%  substitution of Table 3.
% =====================================================================
\subsection{Convergence of EF21-MuonUSign and EF21-MuonSign}\label{app:hyp}

We do not analyze the Error-Feedback methods from scratch. Instead we show that EF21-MuonUSign and EF21-MuonSign are \emph{exact instances} of EF21-Muon \citep[Algorithm~3]{gruntkowska2025error}, whose convergence is already established in the layer-wise, stochastic, federated setting. The whole task is to check that our algorithms meet the framework's hypotheses; the guarantees then transfer verbatim. Only one of these checks is non-trivial -- that the scaled sign is a valid compressor (Lemma~\ref{lem:signcontr}); the rest is bookkeeping. We keep the reduction in our own notation and defer the single change of variables to Table~\ref{tab:dict}.

Throughout, the model is the tuple of matrix layers $\mathbf{X} = [\mathbf{X}_1,\dots,\mathbf{X}_p]$, $\mathbf{X}_i \in \mathbb{R}^{m_i\times n_i}$, with $d_i := m_i n_i$, $r_i := \min(m_i,n_i)$, $d_{\max} := \max_i d_i$; a second subscript selects a layer ($\mathbf{X}_{t,i}$, $\mathbf{G}_{t,i}$). We use $\|\mathbf{Y}\|_{2\to2} \le \|\mathbf{Y}\|_F \le \|\mathbf{Y}\|_{*} \le \sqrt{\operatorname{rank}\mathbf{Y}}\,\|\mathbf{Y}\|_F$ between the spectral, Frobenius and nuclear norms.

\subsubsection{Assumptions}

We use the standard assumptions of the framework, specialized to the spectral (Muon) geometry.

\begin{assumption}[\textbf{Lower boundedness}]\label{as:1}
$f(\mathbf{X}) \ge f^*$ for all $\mathbf{X}$, and (only for Corollary~\ref{cor:l0l1}) each local $f_j \ge f_j^*$.
\end{assumption}

\begin{assumption}[\textbf{Layer-wise smoothness}]\label{as:2}
Each layer carries the spectral norm $\|\cdot\|_{2\to 2}$, dual to the nuclear norm $\|\cdot\|_{*}$. For every layer $i$ and all $\mathbf{X},\mathbf{Y}$,
\begin{equation}
\|\nabla_i g(\mathbf{X}) - \nabla_i g(\mathbf{Y})\|_{*} \le L_i^{g}\,\|\mathbf{X}_i - \mathbf{Y}_i\|_{2\to 2},
\end{equation}
for $g = f$ (constant $L_i$) and $g = f_j$ (constant $L_{i,j}$); set $\tilde{L}_i^2 := \frac{1}{N}\sum_j L_{i,j}^2$ and $L := \max_i L_i$. The weaker layer-wise $(L^0,L^1)$-smoothness \citep{riabinin2025gluon,pethick2025generalized} replaces $L_i^{g}$ by $L^0 + L^1\|\nabla_i g(\mathbf{X})\|_{*}$; it reduces to the above when $L^1 = 0$.
\end{assumption}

\begin{assumption}[\textbf{Stochastic gradient}]\label{as:3}
Each client's stochastic gradient is unbiased, $\mathbb{E}_{\xi}[\nabla f_j(\mathbf{X};\xi)] = \nabla f_j(\mathbf{X})$, with bounded variance $\mathbb{E}_{\xi}\|\nabla f_j(\mathbf{X};\xi) - \nabla f_j(\mathbf{X})\|_*^2 \le \sigma^2$.
\end{assumption}

These are Assumptions~1--2 and~6--10 of \citet{gruntkowska2025error} with the layer norms taken spectral; the variance bound is (mildly) stronger than theirs, which is stated in the Frobenius norm and implied by ours via $\|\cdot\|_F \le \|\cdot\|_{*}$.

\subsubsection{Main result}

\begin{theorem}[\textbf{Convergence of the EF21 methods}]\label{thm:conv}
Run the federated Algorithm~\ref{alg:fed_serverlmo} with the EF21 uplink and EMA momentum $\mu \in [0,1)$. Then both methods are exact instances of EF21-Muon (Proposition~\ref{prop:instance}), and:
\begin{itemize}
\item[\emph{(i)}] \emph{(smooth; both methods)} Under Assumptions~\ref{as:1}--\ref{as:3}, with the sharp learning rate $\eta_{t,i} = \gamma_i\|\mathbf{G}_{t,i}\|_{*}$ and suitably tuned $(\gamma_i,\mu)$, EF21-MuonUSign and the bidirectional EF21-MuonSign both reach
\begin{equation*}
\tfrac{1}{T}\textstyle\sum_{t<T}\mathbb{E}\|\nabla f(\mathbf{X}_t)\|_*^2 = \mathcal{O}(T^{-1/2})
\end{equation*}
(Corollary~\ref{cor:smooth}).
\item[\emph{(ii)}] \emph{(generalized smooth; uplink only)} Under the weaker $(L^0,L^1)$-smoothness, EF21-MuonUSign with a plain \emph{constant} learning rate reaches $\min_{t\le T}\mathbb{E}\|\nabla f(\mathbf{X}_t)\|_* = \mathcal{O}(T^{-1/4})$ (Corollary~\ref{cor:l0l1}).
\end{itemize}
By contrast, the majority-vote placements SignMuon, MuonUSign, MuonSign (Theorems~\ref{th:1}--\ref{th:3}) and EF21-SignMuon (Theorem~\ref{th:ef_div}) admit no such guarantee.
\end{theorem}

Part~(i) gives the $\mathcal{O}(T^{-1/2})$ rate announced in the main text (via $\min_t \le$ average); the centralized Algorithms~\ref{central_alg_ef}--\ref{central_alg_ud} are the case $N=1$. The rest of this section proves the theorem by verifying the framework's four requirements.

\subsubsection{The three routine checks}

\paragraph{The step is a spectral LMO step.}
Writing $\mathbf{G} = \mathbf{U}\boldsymbol{\Sigma}\mathbf{V}^\top$, the framework's oracle over the spectral ball of radius $\tau$ is $\mathrm{lmo}(\mathbf{G}) = \mathbf{X} - \tau\mathbf{U}\mathbf{V}^\top$. Our server step is exactly this with $\tau = \eta_{t,i}$, since $\mathbf{D}_{t,i} = -A(\mathbf{G}_{t,i}) = \mathbf{U}_{t,i}\mathbf{V}_{t,i}^\top$:
\begin{equation}\label{eq:lmo_step_ours}
\mathbf{X}_{t,i} = \mathbf{X}_{t-1,i} - \eta_{t,i}\,\mathbf{D}_{t,i} = \mathrm{lmo}_{\mathcal{B}(\mathbf{X}_{t-1,i},\eta_{t,i})}(\mathbf{G}_{t,i}).
\end{equation}
The framework's ``sharp'' step $\mathbf{X} - \gamma\,\mathbf{G}^{\sharp}$ uses $\mathbf{G}^{\sharp} = \|\mathbf{G}\|_{*}\mathbf{U}\mathbf{V}^\top$ for the spectral norm; it is recovered by the schedule $\eta_{t,i} = \gamma_i\|\mathbf{G}_{t,i}\|_{*}$ (used in Corollary~\ref{cor:smooth}), while a constant $\eta_{t,i}$ is the plain Muon rate (Corollary~\ref{cor:l0l1}). Both are learning-rate choices, not algorithmic changes. For the spectral norm the norm-equivalence constants are $\underline{\rho}_i = 1$, $\bar{\rho}_i = \sqrt{r_i}$. At rank-deficient $\mathbf{G}$ the LMO is non-unique; any selection works, since the analysis uses only $\langle\mathbf{G},\mathrm{lmo}(\mathbf{G})\rangle = -\tau\|\mathbf{G}\|_{*}$ and $\|\mathrm{lmo}(\mathbf{G})\|_{2\to2}\le\tau$.

\paragraph{Momentum and loop order.}
Algorithm~\ref{alg:fed_serverlmo} uses EMA momentum $\mathbf{M}_t = \mu\mathbf{M}_{t-1} + (1-\mu)\mathbf{G}_t$, which is the framework's momentum with $\beta = 1-\mu$. Its heavy-ball form (Equation~\ref{eq:signa_general}) gives \emph{identical iterates}: with zero initialization one checks by induction that the heavy-ball states are the EMA states scaled by $1/(1-\mu)$, and both $\operatorname{sign}(\cdot)$ and the Muon LMO are invariant under positive scaling, so $\mathbf{X}_t$ is unchanged. (The optional Nesterov branch is not covered; see Remark~\ref{rem:scope}.) The only remaining discrepancy is the order of operations within a round, which is resolved next.

\begin{proposition}[\textbf{Exact instance}]\label{prop:instance}
Fix $\mu$, a learning-rate schedule, and the LMO selection above. Started from $\mathbf{M}_0^{(j)} = \mathbf{g}_0^{(j)} = \mathbf{G}_0 = \mathbf{0}$, $\mathbf{W}_0 = \mathbf{X}_0$, Algorithm~\ref{alg:fed_serverlmo} with the EF21 uplink and either downlink mode ($\mathcal{C}^{\downarrow} \in \{\text{exact},\text{EF21-P}\}$) produces the same trajectory as Algorithm~3 of \citet{gruntkowska2025error} with spectral norms, scaled-sign worker compressors, identity/scaled-sign server compressor, $\beta = 1-\mu$, and radii $t_i^k = \eta_{k,i}$ -- up to a one-step index shift $X^{t+1} = \mathbf{X}_t$.
\end{proposition}
\paragraph{Proof} The two loops differ only in where the round is cut: \citet{gruntkowska2025error} order it \emph{step $\to$ downlink $\to$ gradient $\to$ uplink}, we order it \emph{downlink $\to$ gradient $\to$ uplink $\to$ step}. Their iteration $k=0$ is vacuous under our initialization: $\mathbf{G}_0 = \mathbf{0}$ gives $X^1 = X^0 = \mathbf{X}_0$ and a zero downlink residual, so $W^1 = \mathbf{X}_0$. Thereafter each of their iterations $k = t$ performs our round $t$ verbatim -- same momentum, same compressed residual added to the estimator, same LMO step \eqref{eq:lmo_step_ours} -- giving $X^{t+1} = \mathbf{X}_t$, $W^{t+1} = \mathbf{W}_t$ by induction. Hence running their method for $K = T+1$ iterations yields the iterates $\{\mathbf{X}_0,\mathbf{X}_0,\mathbf{X}_1,\dots,\mathbf{X}_{T-1}\}$, and any average or minimum over them equals ours up to one duplicated (nonnegative) term. $\blacksquare$

Zero initialization also makes the framework's initial-error constant explicit: the $\|\mathbf{M}_0 - \mathbf{g}_0\|$ term vanishes and the gradient-deviation terms become $\|\nabla_i f(\mathbf{X}_0)\|$, so $\Psi^0$ is a finite constant.

\subsubsection{The one real check: the scaled sign is a compressor}

The framework needs every transmitted message to come from a \emph{contractive compressor}: a map $\mathcal{C}$ with $\mathbb{E}\|\mathcal{C}(\mathbf{Y})-\mathbf{Y}\|^2 \le (1-\alpha)\|\mathbf{Y}\|^2$ for some $\alpha \in (0,1]$ \citep[Def.~1]{gruntkowska2025error}. This is exactly the property that fails for a bare sign and holds once it is scaled.

\begin{lemma}[\textbf{Contractivity of the scaled sign}]\label{lem:signcontr}
For every $\mathbf{Y} \in \mathbb{R}^{m\times n}$ ($d = mn$), the scaled sign $\mathcal{C}(\mathbf{Y}) = \operatorname{mean}(|\mathbf{Y}|)\operatorname{sign}(\mathbf{Y})$ satisfies
\begin{multline}\label{eq:sign_identity}
\|\mathcal{C}(\mathbf{Y})-\mathbf{Y}\|_F^2 = \|\mathbf{Y}\|_F^2 - \tfrac{2}{d}\|\mathbf{Y}\|_1^2 + \tfrac{\|\mathbf{Y}\|_0}{d^2}\|\mathbf{Y}\|_1^2 \\
\le \bigl(1-\tfrac{1}{d}\bigr)\|\mathbf{Y}\|_F^2,
\end{multline}
so $\mathcal{C}$ is contractive in the Euclidean norm with $\alpha = 1/d$. The bound $\alpha = 1/d$ is met only by $1$-sparse inputs; for dense $\mathbf{Y}$ the true contraction is $\alpha(\mathbf{Y}) = \|\mathbf{Y}\|_1^2/(d\|\mathbf{Y}\|_F^2) = \Theta(1)$ (e.g.\ $\to 2/\pi$ for i.i.d.\ Gaussian entries).
\end{lemma}
\paragraph{Proof} With $c := \|\mathbf{Y}\|_1/d$, zero entries contribute nothing and each nonzero entry contributes $(|Y_{kl}|-c)^2$, giving $\|\mathcal{C}(\mathbf{Y})-\mathbf{Y}\|_F^2 = \|\mathbf{Y}\|_F^2 - 2c\|\mathbf{Y}\|_1 + c^2\|\mathbf{Y}\|_0$, i.e.\ \eqref{eq:sign_identity}. Then $\|\mathbf{Y}\|_0 \le d$ and $\|\mathbf{Y}\|_1 \ge \|\mathbf{Y}\|_F$ give the bound; the dense case sets $\|\mathbf{Y}\|_0 = d$, and the Gaussian limit uses $\mathbb{E}|Y_{kl}| = \sqrt{2/\pi}\,\sigma$. $\blacksquare$

\begin{remark}[\textbf{Why the scaling is decisive}]\label{rem:whyscale}
The \emph{unscaled} sign is not contractive for any $\alpha$: as $\|\mathbf{Y}\|_F \to 0$ on a fixed support, $\|\operatorname{sign}(\mathbf{Y})-\mathbf{Y}\|_F \to \sqrt{\|\mathbf{Y}\|_0}$ does not vanish. This is the exact line between our divergent and convergent methods. SignMuon, MuonUSign and MuonSign broadcast unscaled signs of full quantities -- outside every error-feedback theory, and indeed divergent (Theorems~\ref{th:1}--\ref{th:3}). The EF21 variants instead send the \emph{scaled} sign of a \emph{residual}; the one extra scalar $\operatorname{mean}(|\cdot|)$ per layer per round, negligible on top of one bit per parameter, is what makes the compressor contractive.
\end{remark}

Applied per layer, the uplink scaled sign lies in $\mathbb{B}_2(1/d_i) \subseteq \mathbb{B}_2(\alpha)$ with the common parameter
\begin{equation}\label{eq:alpha_def}
\alpha := 1/d_{\max}.
\end{equation}
The two methods differ only on the downlink: EF21-MuonUSign broadcasts the exact model (server compressor $= \mathcal{I}$, $\alpha_P = 1$), while EF21-MuonSign applies the same scaled sign, contributing a second contractive compressor in $\mathbb{B}_2(\alpha)$.

\subsubsection{Transferred guarantees}

All four requirements hold, so the framework's theorems apply through the change of variables of Table~\ref{tab:dict} ($\underline{\rho}_i = 1$, $\bar{\rho}_i = \sqrt{r_i}$, $\alpha_D = \alpha$). We state the resulting rates and stepsize rules; the explicit non-asymptotic bounds are those of \citet{gruntkowska2025error}, Theorem~19 (smooth) and Theorem~24 ($(L^0,L^1)$), evaluated at these constants, and we do not reproduce them.

\begin{table}[!t]
\centering\footnotesize
\setlength{\tabcolsep}{4pt}
\begin{tabular}{@{}lll@{}}
\toprule
 & \textbf{This paper} & \textbf{\citet{gruntkowska2025error}} \\
\midrule
clients / rounds & $N$, $T$ & $n$, $K = T+1$ \\
iterate & $\mathbf{X}_t$ & $X^{t+1}$ \\
momentum & $\mu$ & $1 - \beta$ \\
learning rate & $\eta_{t,i}$ & radius $t_i^{t}$ \\
layer norm / dual & $\|\cdot\|_{2\to2}$, $\|\cdot\|_{*}$ & $\|\cdot\|_{(i)}$, $\|\cdot\|_{(i)\star}$ \\
norm equivalence & $1,\ \sqrt{r_i}$ & $\underline{\rho}_i,\ \bar{\rho}_i$ \\
uplink compr. & scaled sign & $\mathbb{B}_2(\alpha_D)$ \\
downlink compr. & exact / sign & $\mathcal{I}$ / $\mathbb{B}_2(\alpha_P)$ \\
compression $\alpha$ & $1/d_{\max}$ & $\alpha_D$ \\
smoothness & $L_i,\ \tilde{L}_i$ & $L_i^0,\ \tilde{L}_i^0$ \\
\bottomrule
\end{tabular}
\caption{Change of variables from our notation to that of \citet{gruntkowska2025error}. The one-step index shift is Proposition~\ref{prop:instance}.}
\label{tab:dict}
\end{table}

\begin{corollary}[\textbf{Smooth case; both methods}]\label{cor:smooth}
Let Assumptions~\ref{as:1}--\ref{as:3} hold. Run Algorithm~\ref{alg:fed_serverlmo} with the EF21 uplink, EMA momentum, and the sharp learning rate $\eta_{t,i} = \gamma_i\|\mathbf{G}_{t,i}\|_{*}$ with any constants $\gamma_i$ below the per-layer threshold of \citet[Thm.~19]{gruntkowska2025error} under Table~\ref{tab:dict}. Then, tuning $\mu$ as in \citet[Cor.~2]{gruntkowska2025error},
\begin{equation*}
\tfrac{1}{T}\textstyle\sum_{t<T}\mathbb{E}\|\nabla f(\mathbf{X}_t)\|_*^2 = \mathcal{O}(T^{-1/2}).
\end{equation*}
\end{corollary}
\paragraph{Proof} By Proposition~\ref{prop:instance} the run is an instance of Algorithm~3 for $K = T+1$ iterations, and by Lemma~\ref{lem:signcontr} its compressors satisfy $\alpha_D = \alpha$ (uplink) and $\alpha_P \in \{1,\alpha\}$ (downlink; EF21-MuonSign's Euclidean downlink is admissible by their Remark~23, which adds a factor $\bar{\rho}_i^2 = r_i$ to the primal terms of the threshold). Substituting $\underline{\rho}_i = 1$, $\bar{\rho}_i = \sqrt{r_i}$ into \citet[Thm.~19]{gruntkowska2025error} gives the stepsize threshold and the $\mathcal{O}(T^{-1/2})$ rate; the duplicated $\mathbf{X}_0$-term contributes only a $\frac{T+1}{T}$ factor. $\blacksquare$

The threshold shrinks polynomially in the layer dimension $d_i$ (through $\bar{\rho}_i^2 = r_i$ and the uplink $\alpha = 1/d_{\max}$), and for EF21-MuonSign carries an extra $r_i$ from the compressed downlink -- the price of the spectral geometry and of one-bit messages (Remark~\ref{rem:dimension}). For the bidirectional method this penalty is structural: the scaled sign cannot satisfy the layer-norm contractivity that the unpenalized branch requires, for any layer of rank above one (Remark~\ref{rem:normmismatch}).

\begin{corollary}[\textbf{$(L^0,L^1)$ case; EF21-MuonUSign}]\label{cor:l0l1}
Let Assumption~\ref{as:2} hold in its $(L^0,L^1)$ form. Run EF21-MuonUSign (exact downlink) with $S := T+2$, momentum $\mu = 1 - S^{-1/2}$, and the \emph{constant} per-layer rate $\eta_{t,i} \equiv \eta_i/S^{3/4}$ for any $\eta_i$ small enough that $\eta_i \le 1$ and $\eta_i^2\sqrt{r_i}\,(L^1_{i,\max})^2 = \mathcal{O}(1)$ (the explicit bound is \citet[Thm.~24]{gruntkowska2025error}). Then
\begin{equation*}
\min_{0\le t\le T}\mathbb{E}\|\nabla f(\mathbf{X}_t)\|_* = \mathcal{O}(T^{-1/4}).
\end{equation*}
\end{corollary}
\paragraph{Proof} Their Theorem~24 requires the identity server compressor, which is exactly the exact downlink of EF21-MuonUSign (hence EF21-MuonSign is excluded here; Remark~\ref{rem:scope}). Applying it with $K = T+1$, $\beta = S^{-1/2}$ and the dictionary substitution yields the stated schedule and rate. $\blacksquare$

\begin{remark}[\textbf{The variant we actually run}]\label{rem:const_lr_smooth}
Under plain smoothness ($L^1 = 0$) the constraint in Corollary~\ref{cor:l0l1} collapses to $\eta_i \le 1$. So the \emph{constant-learning-rate} EF21-MuonUSign used in our experiments -- with no gradient-norm-dependent scaling -- already converges at rate $\mathcal{O}(T^{-1/4})$; the sharp schedule of Corollary~\ref{cor:smooth} is needed only to obtain the faster $\mathcal{O}(T^{-1/2})$ rate in squared norm.
\end{remark}

\subsubsection{Scope and consequences}

\begin{remark}[\textbf{What the reduction does and does not cover}]\label{rem:scope}
It covers the two Error-Feedback methods and no others, matching our negative results. \emph{(i)}~The majority-vote methods send unscaled signs and fall outside the framework (Remark~\ref{rem:whyscale}); Theorems~\ref{th:1}--\ref{th:3} confirm they diverge. \emph{(ii)}~EF21-SignMuon compresses the LMO \emph{output} rather than the gradient, so it tracks the non-Lipschitz polar factor, the framework's momentum-tracking step breaks, and it diverges for every $L$-based step size (Theorem~\ref{th:ef_div}). \emph{(iii)}~EF21-MuonSign gets the smooth guarantee (Corollary~\ref{cor:smooth}) but not the $(L^0,L^1)$ one, since the generalized-smooth theory of \citet{gruntkowska2025error} assumes an uncompressed downlink. \emph{(iv)}~The guarantees are for the default EMA branch; the Nesterov branch is not analyzed.
\end{remark}

\begin{remark}[\textbf{The price of one bit}]\label{rem:dimension}
The uplink parameter $\alpha = 1/d_{\max}$ is what makes the stepsize threshold of Corollary~\ref{cor:smooth} shrink with the dimension -- the worst-case cost of one-bit messages, of the same nature as for Top-$1$ compressors in Euclidean EF21 \citep{richtarik2021ef21}. \emph{On the uplink} it is a worst case only: it needs $1$-sparse residuals, whereas the momentum residual is dense with $\alpha(\boldsymbol{\Delta}) = \Theta(1)$ (Lemma~\ref{lem:signcontr}), and the framework admits iteration-dependent $\alpha$ \citep[Rem.~12]{gruntkowska2025error}, recovering the constants realized in practice.

\emph{The downlink of EF21-MuonSign does not inherit this reassurance.} Its residual is not a difference of gradients refreshed by fresh stochasticity each round, but the output of the compressor's own recursion
\begin{equation*}
\boldsymbol{\Delta}_{t+1}^{\downarrow} = \boldsymbol{\Delta}_t^{\downarrow} - \eta_t \mathbf{D}_t - \operatorname{mean}|\boldsymbol{\Delta}_t^{\downarrow}|\operatorname{sign}(\boldsymbol{\Delta}_t^{\downarrow}),
\end{equation*}
in which every coordinate is corrected by the \emph{same} amount. Coordinates driven persistently harder than that mean are never caught while the bulk holds the mean down, so the residual concentrates -- and concentration is what drives $\alpha$ towards its $1/d$ floor. This positive feedback is absent from the uplink, and we observe it: $\alpha(\boldsymbol{\Delta}^{\downarrow})$ falls to $1.2\times10^{-4}$ on the layers built from a zero initialization, four orders of magnitude below its own uplink on the same layers (Section~\ref{sec:exp_nanogpt}). For the bidirectional method the dimension dependence of Corollary~\ref{cor:smooth} is therefore realizable, not pessimistic.
\end{remark}

\begin{remark}[\textbf{Why the sharp route is closed}]\label{rem:normmismatch}
Corollary~\ref{cor:smooth} reaches EF21-MuonSign only through the $\bar{\rho}_i^2$-penalized branch of \citet[Rem.~23]{gruntkowska2025error}, and this is forced rather than convenient. The sharp branch asks for a server compressor in $\mathbb{B}(\alpha_P)$, contractive in the \emph{layer} norm; by the norm-equivalence argument of \citet[App.~C]{gruntkowska2025error} a compressor in $\mathbb{B}_2(\alpha)$ lies there only if $\alpha > 1 - \underline{\rho}_i^2/\bar{\rho}_i^2 = 1 - 1/r_i$. The scaled sign has $\alpha \le 1$, with equality only on the trivial input, so for any layer of rank $r_i>1$ it can \emph{never} qualify -- for our $768\times3072$ matrices the bar is $\alpha > 0.9987$ against an isotropic $2/\pi$. The factor $\bar{\rho}_i^2 = r_i$ is thus not removable by a sharper analysis of this compressor, only by changing it. Two admissible replacements are cheap: random dropout ($\mathcal{C}(\mathbf{Y}) = \mathbf{Y}$ with probability $p$, else $\mathbf{0}$) lies in $\mathbb{B}(p)$ for \emph{every} norm, and the Top-$K$ SVD compressor lies in $\mathbb{B}(1-\sigma_{K+1}^2/\sigma_1^2)$ for the spectral norm \citep[App.~C]{gruntkowska2025error}. Both restore the sharp branch at a comparable budget, giving up the strictly one-bit downlink.
\end{remark}

\begin{remark}[\textbf{From Muon to Gluon}]\label{rem:gluon}
Only the constants $(\underline{\rho}_i,\bar{\rho}_i) = (1,\sqrt{r_i})$ are spectral-norm-specific: Lemma~\ref{lem:signcontr} is a Euclidean statement and the framework's theorems hold for arbitrary layer norms. Replacing the Muon LMO with any per-layer norm -- the Gluon setting \citep{riabinin2025gluon}, e.g.\ Scion's $\ell_1\to\ell_\infty$ geometry for embeddings \citep{pethick2025training} -- gives \emph{EF21-GluonUSign} and \emph{EF21-GluonSign}, for which both corollaries hold with the corresponding $(\underline{\rho}_i,\bar{\rho}_i)$. Convergence is thus a property of error feedback plus scaled-sign compression, not of the spectral geometry.
\end{remark}

We assume the exact spectral LMO, as do all analyses of Muon-type methods \citep{li2025note,kovalev2025understanding,riabinin2025gluon,gruntkowska2025error}; in practice it is the five-step Newton--Schulz approximation of Algorithm~\ref{alg:muon_lmo} \citep{amsel2025polar}, and vector parameters and the last layer are trained with AdamW, as is standard for Muon \citep{jordan2024muon}. Finally, we note that although SignMuon carries no worst-case guarantee, empirically it stays within $1$--$1.5\%$ of full-precision Muon while cutting uplink communication $\sim32\times$ and beating SignSGD, in both the centralized and federated regimes.
